\section{Effective systems recovering the entire base field}
\label{sec:effective}

The exact presentation in the denominator theorem identifies complete
section spaces on the total space and the base. We combine this
identification with effective birationality on a base model.

\begin{proposition}\label{prop:effective-base}
Under the hypotheses of Theorem~\ref{thm:relative-denominators}, suppose in
addition that $D$ is big.  There is a positive integer $m=m(d,\Phi)$ such
that, for every positive multiple $l$ of $m$, the complete divisorial
system
\[
 \left|\left\lfloor l(K_X+B)\right\rfloor\right|
\]
is nonempty and generates precisely the subfield $\C(Z)\subset\C(X)$.
Its rational map is therefore birationally equivalent to $f$ and to the
Iitaka fibration of $K_X+B$.
\end{proposition}

\begin{proof}
Put $L=K_X+B$ and $K=\C(Z)$, and use the exact presentation
\eqref{eq:relative-exact-presentation}. For every positive integer $l$
divisible by $p_0$, we first prove the equality of subspaces of $\C(X)$
\begin{equation}\label{eq:relative-section-spaces}
 H^0\!\left(X,\OO_X(\lfloor lL\rfloor)\right)
   =\psi^{l/p_0}f^*
      H^0\!\left(Z,\OO_Z(\lfloor lD_Z\rfloor)\right).
\end{equation}
Here $f^*$ on the right means pullback of rational functions. For
$v\in K^*$, the exact divisor identity gives
\begin{equation}\label{eq:relative-section-divisors}
 lL+\Div_X\!\left(\psi^{l/p_0}f^*v\right)
    =f^*\!\left(lD_Z+\Div_Z(v)\right).
\end{equation}
If $v$ is a section of the rounded system on $Z$, then
$lD_Z+\Div_Z(v)\geq0$ by the section convention of
Section~\ref{sec:preliminaries}. This divisor is rational Cartier, so
local effective Cartier multiples show that its pullback is effective.
Equation~\eqref{eq:relative-section-divisors} gives the inclusion from
right to left in \eqref{eq:relative-section-spaces}.

Conversely, let $u$ be a nonzero section on the left, and let $X_\eta$
be the generic fibre over $K$. The exact presentation rewrites the
section inequality as
$\Div_X(u/\psi^{l/p_0})+lf^*D_Z\geq0$. Its horizontal valuations give
\[
 \Div_{X_\eta}\!\left(u/\psi^{l/p_0}\right)\geq0.
\]
The generic fibre is normal, so this quotient is regular there. The
contraction condition $f_*\OO_X=\OO_Z$ gives
$H^0(X_\eta,\OO_{X_\eta})=K$. Hence
$u=\psi^{l/p_0}f^*v$ for some $v\in K^*$. By
\eqref{eq:relative-section-divisors}, the source section inequality is
$f^*G\geq0$, where $G=lD_Z+\Div_Z(v)$ is rational Cartier.
Effectivity of $G$ descends to $Z$: near the generic point of any prime
$P\subset Z$, write a Cartier multiple of $G$ as $aP$, with $P$
locally principal. The local Cartier pullback of $P$ has a prime
component dominating $P$, with positive multiplicity $e$, because
$f$ is surjective. Its coefficient in the pullback of the chosen
multiple of $G$ is $ae$. Thus $f^*G\geq0$ forces $a\geq0$.
Doing this for every $P$ proves $G\geq0$, so $v$ is a section on $Z$.
This proves \eqref{eq:relative-section-spaces}. Since $p_0$ clears
$\Phi$, the divisor $lL$ is integral Weil. The argument applies to its
reflexive sheaf even when $lL$ is not Cartier.

We now produce birational sections on the base. Choose a smooth
projective determination $q:W\to Z$ of the nef data which also resolves
the boundary, and write
\[
 K_W+B_W^{\mathrm{cr}}+M_W=q^*D_Z.
\]
Let $A$ be the strict transform of $B_Z$ plus the reduced exceptional
divisor. On nonexceptional primes $A$ agrees with $B_W^{\mathrm{cr}}$;
on exceptional primes its coefficient is one, at least the coefficient
of $B_W^{\mathrm{cr}}$ by generalized log canonicity. Thus
\begin{equation}\label{eq:relative-big-model}
 K_W+A+M_W=q^*D_Z+E_W,
 \qquad E_W=A-B_W^{\mathrm{cr}}\geq0
       \text{ exceptional over }Z.
\end{equation}
Put $\Lambda=\mathcal B\cup\{1\}$. The pair $(W,A)$ is an ordinary
log smooth lc pair with coefficients in the fixed DCC set $\Lambda$,
and $pM_W$ is nef Cartier. Since $D$ is big and $D_Z\sim_\Q D$, the
adjoint on the left of \eqref{eq:relative-big-model} is big.

For each $1\leq k\leq d$, let $b_k=b(\Lambda,k,p)$ be the integer in
Birkar--Zhang's polarized effective birationality theorem
\cite[Theorem~1.3]{BirkarZhang2016}. Its conclusion is that
\[
 \left|\left\lfloor l(K_W+A+M_W)\right\rfloor\right|
\]
is birational for every positive $l$ divisible by $b_k$ when
$\dim W=k$; the floor is taken on the whole adjoint. Set
\[
 m=\operatorname{lcm}(p_0,b_1,\ldots,b_d).
\]
This integer depends only on $d$ and $\Phi$. Fix any positive multiple
$l$ of $m$ and take $k=\dim Z$. For a rational section
$v\in\C(W)=K$ of the displayed birational system, pushforward of its
divisor inequality using \eqref{eq:relative-big-model} gives
\[
 \Div_Z(v)+lD_Z\geq0.
\]
Thus these sections belong to the complete rounded system on $Z$,
and their ratios generate $K$ because the system on $W$ is birational.
The complete system on $Z$ is therefore nonempty and generates $K$.
Equation~\eqref{eq:relative-section-spaces} transfers precisely these
ratios to $X$, since the common factor $\psi^{l/p_0}$ cancels. The
complete source system is nonempty and generates the embedded field
$K\subset\C(X)$ for every such $l$.

The function field of the image of a system's rational map is generated
by its section ratios, so each of these maps is birationally equivalent
to $f$. To compare with the Iitaka fibration, choose for the individual
pair an integer $r>0$ such that $rL$ is Cartier. Degrees divisible by
$\operatorname{lcm}(m,r)$ form a cofinal divisible subsequence of the
Cartier section series of $L$. Such a subsequence has the same field
of section ratios: given $u/v$ in one degree, choose $a\geq1$ so that
the multiplied degree lies in the subsequence and write
$u/v=(uv^{a-1})/v^a$. The field defining the Iitaka
fibration is therefore $K$, and
$\kappa(L)=\operatorname{trdeg}_\C K=\dim Z$. This identifies $f$
with the Iitaka fibration up to birational equivalence. The individual
index $r$ is used only for this comparison and does not enter $m$.
\end{proof}
